% Part: normal-modal-logic
% Chapter: axioms-systems
% Section: introduction

\documentclass[../../../include/open-logic-section]{subfiles}

\begin{document}

\olfileid{nml}{axs}{int}

\olsection{پېژندګلو}

د بنسټيزې مودالي ژبې معناپوهنه مو د مودالي مدلونو په واسطه
لرو، او دا مفهوم هم لرو چې !!a{formula} معتبر وي، يعنې په
ټولو مدلونو کښې په ټولو نړۍو کښې رښتيا وي؛ يا د مدلونو د
کوم ټولګي يا د چوکاټونو په نسبت معتبر وي، يعنې د ټولګي
په ټولو مدلونو کښې يا پر چوکاټ ولاړو ټولو مدلونو کښې
په ټولو نړۍو کښې رښتيا وي۔ منطق عموماً د اعتبار دا
معنايي ځانګړنې د ثبوت‌تيورۍ له !!{derivability} سره تړي۔
موخه دا ده چې په يوه نظام کښې !!{derivability} داسې
تعريف کړو چې !!a{formula} هله او يوازې هله
!!{derivable} وي چې معتبر وي۔

د !!{derivation} تر ټولو ساده او له تاريخي پلوه تر ټولو
زاړۀ نظامونه د هېلبرټ په ډول يا بديهي !!{derivation}
نظامونه دي۔ د ډېرو مودالي منطقونو دپاره د هېلبرټ په ډول
!!{derivation} نظامونه جوړول نسبتاً اسانه دي: د ميتا-تيورۍ
د مطالعې د څيزونو په توګه ساده دي (د بېلګې په توګه، د
سموالي او بشپړتيا د ثبوت دپاره)۔ خو د طبيعي استنتاج د
نظامونو په پرتله په هغو کښې د !!{formula}s ثبوت ډېر سخت دے۔

د هېلبرټ په ډول د !!{derivation} په نظامونو کښې، د يوه
!!a{formula} يو !!a{derivation} د !!{formula}s هغه لړۍ ده
چې له ځينو بديهي اصلونو څخه، د څو استنتاجي قاعدو په
کارولو، تر اړوند !!{formula} رسېږي۔ ځکه چې غواړو د
!!{derivation} نظام له معناپوهنې سره برابر وي، بايد دا
تضمين کړو چې ټول !!{derivable} فارمولونه په ټولو مدلونو
کښې رښتيا وي (يا په هغو ټولو مدلونو کښې رښتيا وي چې
ټول بديهي اصلونه پکې رښتيا وي)۔ لومړے به د مودالي
منطقونو هغه خاصيتونه بېل کړو چې د دې کار دپاره اړين
دي: «عادي» مودالي منطقونه۔ د عادي مودالي منطقونو دپاره
يوازې دوه استنتاجي قاعدې فرضول غواړي: مودوس پوننس
او د لزوم ورکولو قاعده۔ د بديهي اصلونو په توګه د ټولو
تاوتولوژيو (د ځايناستۍ) ټول مثالونه اخلو، او د اړوند
مودالي منطق له مخې يو شمېر مودالي بديهي اصلونه هم۔
که يوازې د ټولو مدلونو ټولګي ته هم ګورو، بيا هم بايد
د~$\Ax{K}$\iftag{notprvDiamond}{}{ او~$\Ax{Dual}$}
د ځايناستۍ ټول مثالونه بديهي اصلونه وشمېرو۔ همدا
د لږ تر لږه عادي مودالي منطق~$\Log K$ بنسټ ږدي۔

\begin{defn}
د \emph{مودوس پوننس} قاعده دا استنتاجي طرحه ده:
\begin{prooftree}
\AxiomC{$!A$}
\AxiomC{$!A \lif !B$}
\RightLabel{\MP}
\BinaryInfC{$!B$}
\end{prooftree}
وايو چې !!a{formula}~$!B$ د مودوس پوننس له مخې
له !!{formula}s $!A$، $!C$ څخه \emph{راځي} هله
او يوازې هله چې $!C \ident !A \lif !B$ وي۔
\end{defn}

\begin{defn}
د \emph{لزوم ورکولو} قاعده دا استنتاجي طرحه ده:
\begin{prooftree}
\AxiomC{$!A$}
\RightLabel{\Nec}
\UnaryInfC{$\Box !A$}
\end{prooftree}
وايو چې !!{formula}~$!B$ د لزوم ورکولو له مخې
له !!{formula}s $!A$ څخه راځي هله او يوازې هله
چې $!B \ident \Box !A$ وي۔
\end{defn}

\begin{defn}
د بديهي اصلونو له سټ~$\Sigma$ څخه يو \emph{!!{derivation}}
د !!{formula}s $!B_1$، $!B_2$، \dots، $!B_n$ داسې لړۍ
ده چې هر~$!B_i$ ئې يا
\begin{enumerate}
\item د يوې تاوتولوژۍ د ځايناستۍ مثال وي، يا
\item په~$\Sigma$ کښې د !!a{formula} د ځايناستۍ مثال وي، يا
\item له دوو !!{formula}s $!B_j$، $!B_k$ څخه، چې $j$، $k < i$ وي،
  د مودوس پوننس له مخې راځي، يا
\item له !!a{formula}~$!B_j$ څخه، چې $j < i$ وي،
  د لزوم ورکولو له مخې راځي۔
\end{enumerate}
که داسې !!a{derivation} وي چې $!B_n \ident !A$، نو وايو
چې~$!A$ \emph{له~$\Sigma$ څخه !!{derivable}} دے؛ په نښو
$\Sigma \Proves !A$۔
\end{defn}

له دې تعريف سره به څرګنده شي چې د ټولو !!{derivable}
فارمولونو سټ يو عادي مودالي منطق جوړوي، او هر
!!{derivable} !!{formula} په هر هغه مدل کښې رښتيا وي
چې هر بديهي اصل پکې رښتيا وي۔ د !!{derivation}s دې
خاصيت ته \emph{سموالے} وايي۔ سرچپه لورې، يعنې
\emph{بشپړتيا}، ثابتول ګران دي۔

\end{document}
